\documentclass{rvcepaper}
\begin{document}

% ---------- HS351TA Entrepreneurship & IPR
\begin{iempaper}
\paperinfo{dept={INDUSTRIAL ENGINEERING AND MANAGEMENT},date={7\textsuperscript{th} Nov 2025},exam={CIE - I},maxmarks={10 + 50},sem={V},prog={UG},duration={30 + 90 Min},
 title={Entrepreneurship \&Intellectual property Rights},code={HS351TA}}
\iemhead
\begin{cietable}{M}{BT}{CO}
\qpart{Part -- A}
\q{1.}{Define Eco labelling}{2}{2}{4}
\q{2.}{List any two characteristics of a successful entrepreneur?}{2}{2}{1}
\q{3.}{Define Compulsory licensing.}{2}{2}{4}
\q{4.}{List any two different types of Entrepreneurs.}{2}{2}{1}
\q{5.}{Rahul, a software engineer, notices that many small businesses struggle with online payments. He develops a mobile app to simplify digital transactions and starts his own company to offer this solution. Based on the above scenario, identify the type of entrepreneur Rahul represents and explain why.}{2}{3}{1}
\qpart{Part -- B}
\q{1}{Define a patent and explain its salient features that make it a valuable legal protection for inventors. Analyze the key factors determining the patentability of an invention.}{10}{3}{4}
\q{2}{Explain the concepts of Registrable and Non-Registrable Trademarks under the Trade Marks Act, 1999. Discuss the procedure for registration of a Trademark under the ACT.}{10}{3}{4}
\q{3.}{A biotechnology company, Bio Nova Pvt. Ltd., holds a patent for a unique process of producing eco-friendly bio-plastic. A rival firm, Green Poly Industries, begins manufacturing and selling a similar bio-plastic using a slightly modified process without obtaining permission from Bio Nova. Around the same time, Bio Nova plans to expand its patent portfolio and conduct a patent search for new innovations while also seeking investors to commercialize and value its intellectual property assets.\newline
Based on the above scenario, analyze the following issues:
\begin{num}
\item What constitutes patent infringement in this case, and what legal remedies are available to Bio Nova under the Patents Act, 1970?
\item How should Bio Nova conduct an effective patent search and prepare a patent draft to protect its future innovations?
\item Discuss the process and importance of commercialization and valuation of Bio Nova's intellectual property in attracting investors.
\end{num}}{10}{4}{5}
\q{4}{Ms. \textbf{Ananya Rao}, a young engineer working in a technology firm, is passionate about developing sustainable energy solutions. While her company supports innovation internally, she is uncertain whether to continue as an intrapreneur within the organization or to start her own entrepreneurial venture focused on solar-powered charging systems for rural areas. She discusses her ideas with a team of colleagues who share her vision, and together they plan to explore market opportunities and assess risks. However, her family believes entrepreneurship is too risky and that only people with large capital or business backgrounds can succeed.\newline
Based on the above scenario, analyze the following issues:
\begin{num}
\item Differentiate between an entrepreneur and an intrapreneur in the context of Ananya's situation.
\item Discuss the emerging trends in entrepreneurship that could influence Ananya's decision to start a new venture.
\item Evaluate the characteristics of a successful entrepreneur that Ananya and her team should develop to ensure business success.
\item Identify and explain the myths about entrepreneurship reflected in the case and how they can be overcome.
\item Explain the role of entrepreneurial teams in fostering innovation and sustaining growth in startups like Ananya's.
\end{num}}{10}{4}{1}
\q{5}{Explain the importance of entrepreneurship in promoting engineering innovation and fostering economic growth. Discuss how entrepreneurial initiatives among engineers contribute to technological advancement, job creation, and sustainable development, with suitable examples.}{10}{3}{1}
\end{cietable}
\end{iempaper}

% ---------- IM352IA Operations Management
\begin{iempaper}
\paperinfo{dept={INDUSTRIAL ENGINEERING \& MANAGEMENT},date={06/11/2025},exam={CIE-- I},maxmarks={10 + 50},sem={V},prog={UG},duration={30 + 90 Min},
 title={Operations Management},code={IM352IA}}
\iemhead
\begin{cietable}{M}{BT}{CO}
\qpart{Part -- A}
\q{1.}{When Amazon decides where to locate its fulfillment centers, it is making a \blank[1.1cm] level decision in operations management.}{1}{2}{2}
\q{2.}{Strategies for change include process improvement, process redesign, and process \blank[1.1cm].}{1}{1}{1}
\q{3.}{Strategic fit refers to the alignment between business strategy and \blank[1.2cm] strategy.}{1}{2}{2}
\q{4.}{In manufacturing, the four process structures are job process, batch process, line process, and \blank[1.2cm] process.}{1}{1}{1}
\q{5.}{Decision-making models in operations include deterministic, probabilistic, and \blank[1.2cm] models.}{1}{1}{3}
\q{6.}{The Fourth Industrial Revolution is driven by technologies such as IoT, AI, and \blank[1.1cm].}{1}{1}{1}
\q{7.}{When a factory balances machine utilization and customer delivery times, it is managing the \blank[1.2cm] responsiveness trade-off.}{1}{2}{2}
\q{8.}{Walmart's use of RFID and data analytics to track inventory is an example of using \blank[1.2cm] as a driver of supply chain performance.}{1}{2}{3}
\q{9.}{In a make-to-order process, production begins only after receiving a \blank[1.2cm] from the customer.}{1}{1}{1}
\q{10.}{Zara's ability to rapidly design and distribute new fashions demonstrates the competitive priority of \blank[1.2cm].}{1}{2}{2}
\qpart{Part -- B}
\q{1}{Define operations management. Explain the key terms involved in the definition of operations management with a schematic model of a production system.}{10}{2}{1}
\q{2}{With suitable example, distinguish between manufacturing and service processes.}{10}{2}{1}
\q{3}{A retailer must decide whether to build a small or a large facility at a new location. Demand at the location can be either low or high, with probabilities estimated to be 0.4 and 0.6, respectively. If a small facility is built and demand proves to be high, the manager may choose not to expand (payoff = \$223,000) or to expand (payoff = \$270,000). If a small facility is built and demand is low, there is no reason to expand and the payoff is \$200,000. If a large facility is built and demand proves to be low, the choice is to do nothing (\$40,000) or to stimulate demand through local advertising. The response to advertising may be either modest or sizable, with their probabilities estimated to be 0.3 and 0.7, respectively. If it is modest, the payoff is estimated to be only \$20,000; the payoff grows to \$220,000 if the response is sizable. Finally, if a large facility is built and demand turns out to be high, the payoff is \$800,000. Draw a decision tree. Then analyze it to determine the expected payoff for each decision and event node. Which alternative---building a small facility or building a large facility---has the higher expected payoff?}{16}{4}{3}
\q{4}{A manager is deciding whether to build a small or a large facility. Much depends on the future demand that the facility must serve, and demand may be small or large. The manager knows with certainty the payoffs that will result under each alternative, shown in the following payoff table. What is the best alternative for each decision rule: Maximin. Maximax. Laplace. Minimax Regret.
\qtab{|l|c|c|}{\hline Alternative & \multicolumn{2}{c|}{Possible future demand}\\\cline{2-3} & Low & High\\\hline Small facility & 200 & 270\\\hline Large facility & 160 & 800\\\hline Do nothing & 0 & 0\\\hline}}{14}{4}{3}
\end{cietable}
\end{iempaper}

% ---------- IM353IA Quality Assurance
\begin{iempaper}
\paperinfo{dept={INDUSTRIAL ENGINEERING \& MANAGEMENT},date={7\textsuperscript{th} Nov 2025},exam={CIE-1},maxmarks={10 + 50},sem={V},prog={UG},duration={30 + 90 Min},
 title={Quality Assurance},code={IM353IA}}
\iemhead
\begin{cietable}{M}{BT}{CO}
\qpart{Part -- A}
\q{1.}{List the statistical methods for quality control and improvement.}{2}{1}{01}
\q{2.}{What is the primary objective of ISO 14000 standards?}{2}{2}{01}
\q{3.}{Write the steps followed in constructing the individual chart}{2}{2}{02}
\q{4.}{Differentiate between defect and defective.}{2}{2}{02}
\q{5.}{What is the primary goal of quality assurance?}{2}{2}{02}
\qpart{Part -- B}
\q{1}{``Analyze how Rational Subgrouping can be used as a quality control strategy in manufacturing or service industries to separate common and special causes of variation, thereby improving process stability.''}{10}{4}{3}
\q{2}{``Evaluate the organizational, operational, and strategic benefits of implementing the ISO 9000 quality standard, and explain how it contributes to continuous improvement and competitive advantage.''}{10}{3}{3}
\q{3}{``Analyze how the Eight Dimensions of Quality can be used to evaluate the performance of a consumer product. Examine each dimension's effect on customer satisfaction and illustrate your analysis with a detailed case study from the consumer goods or electronics industry.''}{10}{4}{3}
\q{4}{A manufacturing company measures the diameter (in mm) of a component during production. Data are collected in
\qtab{|c|c|c|c|c|}{\hline Sample & Measurement 1 & Measurement 2 & Measurement 3 & Measurement 4\\\hline 1 & 24.8 & 25.1 & 24.9 & 25.0\\\hline 2 & 25.3 & 25.2 & 25.4 & 25.1\\\hline 3 & 24.7 & 24.9 & 24.8 & 24.6\\\hline 4 & 25.5 & 25.4 & 25.6 & 25.3\\\hline 5 & 24.9 & 25.0 & 25.1 & 24.8\\\hline 6 & 25.2 & 25.3 & 25.1 & 25.0\\\hline}
\begin{ABC}
\item sample calculate the sample mean
\item determine the $\bar{X}$-chart control limits and the R-chart control limits:
\item Interpret the control charts: are the process means and ranges in control? Identify any points or patterns that indicate special-cause variation.
\end{ABC}}{10}{4}{4}
\q{5}{A process is producing shafts with a nominal diameter of 50 mm.\newline
The specification limits are $50\pm0.5$ mm.\newline
From control chart data, the process mean ($X$) is 50.1 mm, and the process standard deviation ($\sigma$) is 0.1 mm.\newline
(a) Calculate the Upper Control Limit (UCL) and Lower Control Limit (LCL) for the process using $\pm3\sigma$ limits.\newline
(b) Compare the control limits with the specification limits and comment on the process capability.}{10}{4}{4}
\end{cietable}
\end{iempaper}

% ---------- IM254TA Financial Accounting & Costing
\begin{iempaper}
\paperinfo{dept={INDUSTRIAL ENGINEERING \& MANAGEMENT},date={07/11/2025},exam={CIE -- I},maxmarks={10 + 50},sem={V},prog={UG},duration={30 + 90 Min},
 title={FINANCIAL ACCOUNTING \& COSTING},code={IM254TA}}
\iemhead
\begin{cietable}{M}{BT}{CO}
\qpart{Part -- A}
\q{1.}{State the primary objective of financial accounting}{1}{1}{1}
\q{2.}{The process of recording a transaction is called \blank[3cm].}{1}{2}{2}
\q{3.}{Showcase your understanding of accounting equation with an example}{2}{2}{2}
\q{4.}{Identify the books which are termed as ``First book of accounting and Principal book of accounting''?}{2}{2}{1}
\q{5.}{State any two advantages of adopting IFRS for Indian engineering firms with global operations.}{2}{2}{3}
\q{6.}{In the era of many startups trying to be unicorns, comment on the ``Conservatism/Prudence concept'' and its importance which has helped traditional businesses prevail in the long run.}{2}{2}{2}
\qpart{Part -- B}
\q{1.}{Review the Golden Rules of Accounting. Illustrate each category with necessary examples.}{10}{2}{1}
\q{2.}{Assume you are starting an E-Commerce website to sell a product of your choice. Identify and list any 10 financial transactions that you might encounter in the first month of starting the business and pass them on to the journal books.}{10}{4}{2}
\q{3.}{The following is a list of various accounts. State the nature of each account (personal, real or nominal) and also which account will be debited and which account will be credited. (\textbf{Use a tabular column to furnish your answer.})\newline
Wages outstanding, Capital introduced, Cartage paid, Prepaid salary, Rent outstanding, Bank overdraft, Discount allowed, Goods returned by buyer, Drawings by proprietor, Commission received.}{10}{3}{3}
\q{4.}{Give the Journal entries for the following transactions:
\qtab{|l|l|l|}{\hline Date & Particulars & \rupee\\\hline April 1 & Business started with cash & 10,000\\\hline April 3 & Deposited in SBI & 6,000\\\hline April 5 & Bought goods from Mahavir & 1,500\\\hline April 9 & Sold goods to Gupta & 650\\\hline April 12 & Paid cash to Mahavir & 990\\\hline April 15 & Cash received from Gupta & 625\\\hline April 20 & Furniture purchased & 300\\\hline April 22 & Drew cash from bank for personal use & 600\\\hline April 30 & Paid rent by cheque & 200\\\hline April 30 & Salary due to clerk & 300\\\hline}
Also prepare the necessary ledger accounts for cash and capital accounts only.}{10}{3}{3}
\q{5.}{Financial literacy is very important for everyone given the road map India has chosen to grow. There are several challenges on the ground also with regard to providing ``Financial Literacy''.\newline
Conceptualize a framework which utilizes students with both engineering \& finance knowledge along with problem solving skills like you to provide ``Financial knowledge'' to students between the age group of 15 to 19. You are allowed to use the same flow which you used for your ``Design Thinking Project''. Comment on Novelty that you want to bring from your POV.}{10}{4}{3}
\end{cietable}
\end{iempaper}

% ---------- IM355TBB Enterprise Information Systems
\begin{iempaper}
\paperinfo{dept={INDUSTRIAL ENGINEERING \& MANAGEMENT},date={10/11/2025},exam={CIE -- I},maxmarks={10 + 50},sem={V},prog={UG},duration={30 + 90 Min},
 title={Enterprise Information Systems},code={IM355TBB},note={Answer all the questions}}
\iemhead
\begin{cietable}{M}{BT}{CO}
\qpart{Part -- A}
\q{1.}{The evolution of planning systems progressed from MRP-I to \blank[1.8cm] and then to ERP.}{1}{1}{CO1}
\q{2.}{The primary mechanism by which Enterprise Information Technology provides value to organizations is through appropriate \blank[3.3cm] improvement.}{1}{1}{CO1}
\q{3.}{Enterprise Systems are noted for providing a tightly integrated solution to an organization's information needs. State two functions that were historically stand-alone applications.}{2}{2}{CO2}
\q{4.}{What is the Role of EIS in business decision-making? State one key function.}{2}{1}{CO1}
\q{5.}{A textile company is implementing a new EIS. Their goal is to maximize the utilization of their collective employee knowledge for project teaming and innovation. Which core EIS component is primarily responsible for managing and exploiting this intellectual capital? Justify your choice.}{2}{3}{CO3}
\q{6.}{What is the Role of EIS in business decision-making? State one key function.}{2}{1}{CO1}
\qpart{Part -- B}
\q{1}{Distinguish between MRP-I and MRP-II. Explain the primary functions of each planning system and illustrate how the evolution from Material Requirements Planning (MRP-I) to Manufacturing Resource Planning (MRP-II) broadened the scope of enterprise planning.}{10}{2}{CO2}
\q{2}{Define and Differentiate the Core Components of EIS. Briefly define ERP, CRM, and SCM, and explain the unique focus of each component within the overall enterprise system framework.}{10}{2}{CO2}
\q{3}{Explain the scope of Digital Transformation in the Context of EIS. Describe what the term means for a modern enterprise and how the adoption of an integrated EIS serves as the foundational enabler for achieving digital transformation goals.}{10}{2}{CO1}
\q{4}{You are advising a high-growth tech startup. The founders are considering which EIS component to prioritize first: CRM for managing investor relations and early customer feedback, or a modular ERP for integrating the finance and HR functions. Argue for the strategic component that should be prioritized to support the startup's growth and eventual Digital Transformation.}{10}{3}{CO3}
\q{5}{What is a Business Process? Provide a comprehensive definition, and explain why identifying and defining these processes is the crucial first step before any automation or reengineering effort.}{10}{2}{CO2}
\end{cietable}
\end{iempaper}

% ---------- IM355TBD Advanced Decision Modelling
\begin{iempaper}
\paperinfo{dept={INDUSTRIAL ENGINEERING \& MANAGEMENT},date={07/11/2025},exam={CIE -- 1},maxmarks={10 + 50},sem={V},prog={UG},duration={30 + 90 Min},
 title={ADVANCED DECISION MODELLING},code={IM355TBD},note={Answer all the Questions.}}
\iemhead
\begin{cietable}{M}{BT}{CO}
\qpart{Part -- A}
\q{1.}{What is balking, reneging and jockeying?}{02}{L2}{CO1}
\q{2.}{The Probability that exactly zero units in the queuing system is \blank[3.5cm]}{02}{L2}{CO1}
\q{3.}{Write the Kendall's notation for the single server queuing system}{02}{L2}{CO1}
\q{4.}{Classify the states of Markov Chain.}{02}{L2}{CO2}
\q{5.}{In an M/M/1 queue, the arrival is Poisson with 4 per minute and the average time spent in system is 5.5 minutes for each customer. The service rate is \blank[1.2cm] per minute}{02}{L3}{CO1}
\qpart{Part -- B}
\q{1}{Customers arrive at a one-window drive according to a Poisson distribution with mean of 10 minutes and service time per customer is exponential with mean of 6 minutes. The space in front of the window can accommodate only three vehicles including the serviced one. Other vehicles have to wait outside this space. Calculate:
\begin{rom}
\item[i.] Probability that an arriving customer can drive directly to the space in front of the window.
\item[ii.] Probability that an arriving customer will have to wait outside the directed space.
\item[iii.] How long an arriving customer is expected to wait before getting the service?
\end{rom}}{10}{L3}{CO1}
\q{2}{At a stamp vending window of a post office 20 customer arrive on an average every 10 min. the vendor serves 5 customers in 2 min. determine:\newline
1. Average number of customers in the system\newline
2. Average waiting time of a customer.\newline
3. Probability of a customer waiting for more than 3 min before being served.\newline
4. Idle time of the vendor in a shift of 8 hours}{14}{L3}{CO1}
\q{3}{Discuss the classification of queuing models and explain how Kendall's notation helps in representing them.}{06}{L3}{CO1}
\q{4}{In a bank there is only one window. A solitary employer performs all the services required and the window remain continuously open from 7am to 1 pm. It was discovered that an average number of clients is 54 during the day and the average service time is 5min / person. Find:
\begin{rom}
\item[i.] Average number of clients in the system
\item[ii.] Average waiting time
\item[iii.] The probability that a client has to spend more than 10 mins in the system
\end{rom}}{10}{L3}{CO1}
\q{5}{A system has three possible states, 0, 1 and 2. Every hour it makes a transition to a different state, which is determined by a coin flip. For example, from state 0, it makes a transition to state 1 or state 2 with probabilities 0.5 and 0.5.\newline
a) Find the transition probability matrix.\newline
b) Find the three-step transition probability matrix.\newline
c) Find the steady-state distribution of the Markov chain.}{10}{L3}{CO2}
\end{cietable}
\end{iempaper}

\end{document}
